\documentclass[final,onefignum,onetabnum]{siamart}

% ------------------------------------------------------------
% Packages (SIAM-approved)
% ------------------------------------------------------------
\usepackage{amsmath, amssymb}
\usepackage{graphicx}

\usepackage{hyperref}

% ------------------------------------------------------------
% Law Environment (SIAM-native alias)
% ------------------------------------------------------------
\newtheorem{law}{Law}
\newtheorem{axiom}{Axiom}
\newtheorem{remark}{Remark}

% ------------------------------------------------------------
% Title and Author
% ------------------------------------------------------------
\title{Supplementary Appendix: The Amsler Surface:\\
The Primitive Computational Power of the Superposed Hypersphere}

\author{
Aaron Schoenman\thanks{Independent Researcher, Columbus, Ohio, USA.
Email: \texttt{yummyinmytummy@naturallybaked.co.site}. ORCID: 0009-0004-8636-0927}
}

\begin{document}
\maketitle

\appendix

\section{Subspace Layer Formation}

The five negative-domain tiers of the primitive hypersphere are shown in Figure~\ref{fig:subspace-tiers}. These curves correspond to the inverted polynomials listed in Appendix~B and illustrate the collapse sequence
from Tier~$-5$ through Tier~$-1$.

\begin{figure}[htbp]
  \centering
  \includegraphics[width=0.85\textwidth]{Figures/SubspaceTiers.png}
  \caption{Figure A.1: Tier curves and inversion boundary at $x=2\pi$.}
  \label{fig:subspace-tiers}
\end{figure}

\section{Inverted Polynomial Expressions}

Inverted polynomial expressions used to generate the subspace curvature graph in Appendix A. Each tier equation represents the reciprocal curvature behavior obtained by inverting the Python radius sweep functions, revealing the subspace structure that emerges as curvature approaches the Amsler asymptote at $x = 2\pi$.

\begin{table}[htbp]
\centering
\caption{Inverted polynomial expressions used to generate the subspace curvature
graph in Appendix~A. Each tier equation represents the reciprocal curvature
behavior obtained by inverting the Python radius sweep functions, revealing the
subspace structure that emerges as curvature approaches the Amsler asymptote at
$x = 2\pi$. Color names correspond to the tier curves shown in Figure~A\textendash1.}
\label{tab:inverted-polynomials}
\begin{tabular}{lll}
\hline
\textbf{Tier} & \textbf{Expression} & \textbf{Color (Figure A\textendash1)} \\
\hline
Base/Seed Tier & $x = 2\pi$ \text{ (Amsler's Asymptote)} & Black \\
Tier $-1$ & $x = 2\pi y^{-1} + 2\pi$ & Red \\
Tier $-2$ & $x = 2\pi y^{-2} + 2\pi$ & Orange \\
Tier $-3$ & $x = 4\pi y^{-3} + 2\pi$ & Green \\
Tier $-4$ & $x = -12\pi y^{-4} + 2\pi$ & Blue \\
Tier $-5$ & $x = 48\pi y^{-5} + 2\pi$ & Purple \\
\hline
\end{tabular}
\end{table}

\section{DC Emergent Amsler Saddle}

\begin{figure}[htbp]
  \centering
  \includegraphics[width=0.85\textwidth]{Figures/DCAmslerSaddle.png}
  \caption{Figure A-2. The subspace tier curves illustrate the full negative-domain structure generated by curvature inversion. Each tier represents a distinct residue of the dimensional ladder, with $T_{-4}$ (black) and $T_{-3}$ (purple) forming the dominant curvature signatures that  survive projection into Euclidean space. As the curves approach the inversion  boundary at $x = 2\pi$, their behavior diverges sharply: $T_{-3}$ rises toward  the Amsler asymptote, while $T_{-4}$ collapses into the reciprocal remainder  that produces the saddle’s asymmetric slopes. The remaining tiers ($-1$, $-2$,  $-5$) contribute progressively smaller residues, forming the nested subspace layers that define the geometry of the negative-curvature domain. Together,  these curves reveal the full structure of inverted curvature broadcast, the  collapse–inversion threshold at $2\pi$, and the projection limits that  generate the Amsler surface as the Tier~$-3$/Tier~$-4$ Euclidean remainder.  Color is required to distinguish the tier curves.}
  \label{fig:dca-amsler-saddle}
\end{figure}

\section{The Dual Hypersphere}

\begin{figure}[htbp]
  \centering
  \includegraphics[width=0.85\textwidth]{Figures/InfinityFootball.jpeg}
  \caption{Figure B\textendash1. The dual hypersphere diagram and the Python generated tier curves together reveal the full orientation mechanics of the  curvature ladder. The broadcast hypersphere occupies the hyper domain above the DC baseline, generating positive curvature through ascent. Its inverted counterpart occupies the subspace domain below the baseline, producing reciprocal curvature through dimensional descent. Euclid’s plane forms the mirror boundary between these domains, enforcing the mandatory $90^\circ$  axis inversion that converts broadcast curvature into its negative-tier  remainder. The Python tier curves quantify this descent: Tier~$-3$ rises toward the Amsler asymptote at $x = 2\pi$, while Tier~$-4$ collapses into the  reciprocal residue that defines the saddle’s asymmetric slopes. The remaining  negative tiers converge toward the subspace point shown in the diagram,  forming nested curvature layers that terminate at the inversion boundary.  Together, the diagram and the tier curves demonstrate that subspace is not an  abstract construct but the geometric and computational footprint of inverted  hyperspherical curvature projected onto the DC baseline. Color is required to  distinguish the domain regions and curvature layers.}
  \label{fig:dual-hypersphere}
\end{figure}

\section{The Tiered Hypersphere Transitional Ladder}

\begin{table}[htbp]
\centering
\caption{Polynomial expressions governing the curvature behavior across the inversion and expansion domains. The negative tiers ($-5$ through $-1$) describe the reciprocal subspace curvature layers generated by dimensional inversion, while the nonnegative tiers represent the forward broadcast expansion of curvature through the Euclidean and hyperspherical domains. Together, these expressions form the complete curvature ladder underlying Dimensional Calculus.}

\label{tab:curvature-ladder}
\begin{tabular}{lll}
\hline
\textbf{Domain} & \textbf{Tier} & \textbf{Polynomial Expression} \\
\hline
Inversion Domain & $-5$ & $48\pi \odot r^{-5}$ \\
                 & $-4$ & $-12\pi \odot r^{-4}$ \\
                 & $-3$ & $4\pi \odot r^{-3}$ \\
                 & $-2$ & $-2\pi \odot r^{-2}$ \\
                 & $-1$ & $2\pi \odot r^{-1}$ \\
Baseline         & $0$  & $2\pi r^{0} = 2\pi \odot$ \\
Expansion Domain & $1$  & $2\pi r^{1}$ \\
                 & $2$  & $\pi r^{2} + C_{1} r^{0}$ \\
                 & $3$  & $\frac{\pi}{3} r^{3} + C_{1} r^{1} + C_{2} r^{0}$ \\
                 & $4$  & $\frac{\pi}{12} r^{4} + \frac{C_{1}}{2} r^{2} + C_{2} r^{1} + C_{3} r^{0}$ \\
                 & $5$  & $\frac{\pi}{60} r^{5} + \frac{C_{1}}{6} r^{3} + \frac{C_{2}}{2} r^{2} + C_{3} r^{1} + C_{4} r^{0}$ \\
\hline
\end{tabular}
\end{table}

The dimensional ladder organizes curvature into a unified sequence of inversion and propagation tiers, each arising from elementary differentiation or integration of the seed curvature term. The negative tiers ($T_{-1}$ through $T_{-5}$) form the reciprocal domain, producing the inverted hypersphere and its nested residues that define subspace geometry. The baseline tier $T_{0}$ marks the DC mirror boundary, where curvature changes sign and the collapse--inversion threshold at $2\pi$ is enforced. Above this boundary, the positive tiers ($T_{1}$ through $T_{5}$) generate the broadcast hypersphere through successive integrations, each tier inheriting constants from the one below it and contributing to the full curvature structure observed in Euclidean space. Together, the inversion and propagation tiers form a continuous curvature ladder whose superposition governs all geometric behavior across hyper, Euclidean, and subspace domains.

\section{Broadcast Curvature Constants/Gravitational Load Equations}

The broadcast curvature constants $C_{n}$ quantify the inherited curvature components that propagate upward through the dimensional ladder. Each constant encodes the residual broadcast from the tier below it, decomposing Euclid’s hypersphere expressions into their radius-governed curvature contributions under Dimensional Calculus. These constants govern the gravitational load across the positive tiers, determining how curvature accumulates, propagates, and superposes as the hypersphere expands through the Euclidean and hyper domains.

\begin{table}[htbp]
\centering
\caption{Broadcast curvature constants $C_{n}$ for the first four tiers of the dimensional ladder. Each polynomial expression represents the inherited curvature broadcast at tier $n$, revealing how Euclid’s hypersphere terms decompose into radius-broadcasted curvature components under Dimensional Calculus.}
\label{tab:broadcast-constants}
\begin{tabular}{ll}
\hline
\textbf{Broadcast Curvature Constant} & \textbf{Polynomial Expression} \\
\hline
$C_{1}$ & $\pi - 2$ \\
$C_{2}$ & $r\{\pi r^{2} - (\pi - 2)\}$ \\
$C_{3}$ & $-r^{2}\left\{\frac{3\pi r^{2} + 2(\pi - 2)}{4}\right\}$ \\
$C_{4}$ & $r^{4}\left\{\frac{(\pi - 2)(-27r^{3} + 45r^{2} + 20r - 30) + (-54r^{3} + 90r^{2} + 60r - 120)}{60}\right\}$ \\
\hline
\end{tabular}
\end{table}

\section{DC's Radius Compared To Euclid's Radius}

Appendix E presents the Python generated curvature deficit graph used to evaluate the computational behavior of the hypersphere. The graph compares the curvature aware radius sweep from Dimensional Calculus with the Euclidean radius function, highlighting the deficit Euclid’s model cannot measure. This visualization serves as the computational counterpart to the analytical results established in Curvature Dynamics and Gibbs: The Hypersphere Equation Is Closed for Renovations (A Fond Farewell to Harmonic Analysis?). The PDE expression traditionally used to model curvature is displayed within the graph’s own frame for direct comparison, illustrating how the hypersphere reduces differential evolution to algebraic evaluation. The accompanying Python code provides full reproducibility and demonstrates the simplification of curvature computation under the DC framework.

\begin{figure}[htbp]
  \centering
  \includegraphics[width=0.85\textwidth]{Figures/DCvEuclidRadius.png}
  \caption{Figure E\textendash1. Curvature deficit visualization comparing the Dimensional Calculus radius sweep with the Euclidean radius function. The shaded region highlights the deficit Euclid’s model cannot measure, and the embedded PDE expression shows the classical curvature formulation for direct  comparison. The graph demonstrates the computational simplification achieved by the hypersphere under DC. Color is required to distinguish the DC radius, Euclidean radius, and curvature-deficit region.}
  \label{fig:dc-vs-euclid-radius}
\end{figure}

\clearpage
\section{Amsler's Surface Does Not Embed; It Becomes The Hypersphere}

To demonstrate how Amsler’s surface becomes the hypersphere under Dimensional Calculus (DC), we must first account for the correct mapping of negative curvature space. In DC, the classical saddle associated with the Amsler surface is not interpreted as an embedded surface at all. Instead, it is recognized as the transitional fold of subspace between Tier~$-3$ and Tier~$-4$ of the curvature ladder. This fold arises directly from the absolute-value structure of the Tier~$-3$ and Tier~$-4$ broadcast expressions, whose opposing curvature directions generate the saddle remainder observed in Euclidean projection.

To expose the curvature deficit left unmeasured by Amsler, we begin by placing the Tier~$-3$ and Tier~$-4$ broadcast expressions into their uninverted DC forms and restricting them to the subspace domain $x<0$. In this domain, the Tier~$-3$ expression produces a divergent curvature contribution, while the Tier~$-4$ expression produces a convergent one. Geometrically, these appear as directional $-\infty$ and $+\infty$ limits relative to the zeroth baseline ($2\pi$). When these two tier contributions are added, their opposing hypersphere infinities cancel, leaving precisely the finite Euclidean remainder that Amsler recorded as the classical saddle. This remainder is the curvature deficit: the portion of the hypersphere broadcast that Euclid cannot embed and Amsler could not measure.

\begin{remark}
We now restrict attention to the subspace domain $x<0$. This restriction is not merely technical; it is geometrically necessary. In the uninverted DC tier expressions, the limit $x\to 0$ corresponds to a blow-up in the curvature broadcast, and allowing $x=0$ would falsely suggest that the transitional fold could coincide with the core hypersphere radius at $2\pi$. In DC, the Tier~$-3$ and Tier~$-4$ fold is strictly a subspace phenomenon: it approaches the zeroth baseline but never attains it. Thus, we enforce $x<0$ to ensure that the Amsler fold remains disjoint from the core hypersphere and is treated correctly as a curvature transition rather than an embedded radius.
\end{remark}

We now write the Tier~$-3$ and Tier~$-4$ broadcast contributions in their uninverted DC forms on the subspace domain $x<0$:

\begin{equation}
    y_{-3}(x) = \frac{4\pi}{x^{3}}, 
    \qquad
    y_{-4}(x) = -\frac{12\pi}{x^{4}},
    \qquad x<0.
\end{equation}

The combined curvature broadcast across the Tier~$-3$/Tier~$-4$ fold is then

\begin{equation}
    y_{\mathrm{fold}}(x)
    = y_{-3}(x) + y_{-4}(x)
    = \frac{4\pi}{x^{3}} - \frac{12\pi}{x^{4}},
    \qquad x<0.
\end{equation}

In the limit as $x\to 0^{-}$, the individual tier contributions exhibit opposing hypersphere infinities relative to the zeroth baseline:

\begin{equation}
    \lim_{x\to 0^{-}} y_{-3}(x) = -\infty,
    \qquad
    \lim_{x\to 0^{-}} y_{-4}(x) = +\infty,
\end{equation}

so that the DC interpretation of the fold is governed by the combined limit

\begin{equation}
    \lim_{x\to 0^{-}} y_{\mathrm{fold}}(x)
    = \lim_{x\to 0^{-}} \big( y_{-3}(x) + y_{-4}(x) \big),
\end{equation}

which defines the curvature deficit: the finite Euclidean remainder of the Tier~$-3$/Tier~$-4$ hypersphere broadcast that Amsler’s saddle records but does not fully resolve.

Now that we have accounted for the curvature that Amsler was unable to measure, DC is free to begin adding the dimensional expansion terms tier by tier. In doing so, Amsler’s surface does not embed into Euclidean space; rather, it integrates into the hypersphere itself. The recovered Tier~$-3$/Tier~$-4$ broadcast provides the missing curvature needed for the saddle to be understood not as an isolated surface, but as the first visible remainder of the full hypersphere transition.

\begin{equation}
    y_{\mathrm{DC}}(x)
    = y_{\mathrm{deficit}}(x)
    + y_{\mathrm{Amsler}}(x)
    + \sum_{k=-2}^{K} T_k(x),
    \qquad x<0,
\end{equation}

\begin{remark}
Throughout this expansion we use the index $k$ to denote the dynamical tier shift in the curvature broadcast. This is distinct from the index $n$ used in the polynomial coefficient set $\{C_n\}$, which enumerates the algebraic limits of the dimensional expansion. The tier index $k$ therefore tracks geometric position within the DC curvature ladder, while the coefficient index $n$ captures the analytic structure of the corresponding expansion term. These indices coincide in ordering but represent different aspects of the DC framework and should not be conflated.
\end{remark}

With all limits defined, we may now demonstrate how the DC reconstruction proceeds. The Amsler saddle provides the Euclidean remainder of the Tier~$-3$/Tier~$-4$ broadcast, while the curvature deficit restores the portion of the hypersphere that Amsler could not measure. DC then adds the remaining dimensional expansion terms tier by tier, beginning at Tier~$-2$ and continuing upward through the curvature ladder. Each tier contributes a broadcast term $T_k(x)$, and each of these terms is represented analytically by a polynomial coefficient $C_n$ drawn from the dimensional expansion set defined within Table~2 and Table~3 from Appendix~F .

Thus the full DC reconstruction of the hypersphere portion associated with the metric length $x$ is given by
\begin{equation}
    y_{\mathrm{DC}}(x)
    = y_{\mathrm{deficit}}(x)
    + y_{\mathrm{Amsler}}(x)
    + \sum_{k=-2}^{K} T_k(x),
    \qquad x<0,
\end{equation}
where each tier contribution satisfies
\begin{equation}
    T_k(x) = C_k(x).
\end{equation}

Substituting the $C_k$ coefficients into the expansion yields a single polynomial whose value at $x$ expresses precisely the portion of the hypersphere broadcast associated with that metric length. In DC, this polynomial is the completed curvature representation: it shows how any segment of the metric, when supplied with the missing broadcast terms, integrates into the hypersphere rather than embedding in Euclidean space.

\subsection{The DC Polynomial as the Curvature Broadcast of the Hypersphere}

In the DC curvature ladder, Tier~$-4$ represents collapsed curvature: a domain with no broadcast, no orientation, and no outward propagation. Curvature exists only as stored geometric memory. Tier~$-3$, immediately above it, retains structural curvature capable of folding and interacting with broadcast. The boundary between these two tiers is therefore the location at which curvature flow reverses direction.

When Tier~$-4$ is actualized in physical spacetime, this boundary corresponds precisely to the black hole event horizon. Outside the horizon, curvature broadcasts outward through the positive tiers; inside, curvature collapses inward through Tier~$-4$. The horizon is the transitional fold between these regimes. In DC, this fold is the same Tier~$-3$/Tier~$-4$ structure whose Euclidean projection Amsler recorded as the classical saddle. Thus the Amsler surface is revealed as the observable shadow of the Tier~$-4$ curvature transition: the geometric signature of the event horizon when projected into Euclidean space.

\begin{remark}
In reconstructing the hypersphere through the DC framework, the obstruction traditionally associated with the Navier--Stokes equations disappears entirely. The apparent ``roughness'' arises only when the flow is forced to diverge from the precise hypersphere geometry. Restoring convergence to the correct curvature broadcast eliminates these pathologies entirely and yields a smooth geometric transition from Amsler's surface into the Tier~4 structure of our physical reality.
\end{remark}

\end{document}